\chapter{Torelli map and universal Jacobian family}
\label{chap:torelli}
% archon:covers MR4276287UniformityInMordellLangForCurves/Torelli.lean

This chapter establishes the Siegel moduli space \(\mathbb{A}_g\) of principally
polarized abelian varieties with level structure, the universal abelian family
over it, and the Torelli morphism \(\tau \colon \mathbb{M}_g \to \mathbb{A}_g\).
These are prerequisites for the ambient setup of \cref{def:ambient_setup} and
the moduli-height comparison of \cref{lem:moduli_height_comparison}.

\section{Preparation for counting points}

\begin{definition}[Siegel moduli space with level structure]
  \label{def:Ag_level}
  \lean{MR4276287UniformityInMordellLangForCurves.Torelli.agLevel}
  % SOURCE: Section 6 of the paper (read from references/MR4276287-mordell-lang-curves-source/SettingUp.tex)
  % SOURCE QUOTE: "Recall from $\mathsection$\ref{subsec:hgtineqnondeg} that $\mathbb{A}_g$ denotes the fine moduli space of principally polarized abelian varieties of dimension $g$, with level-$\ell$-structure. Moreover, $\mathbb{A}_g$ is regular and quasi-projective; see \cite[Theorem~7.9 and below]{MFK:GIT94} or \cite[Theorem 1.9]{OortSteenbrink}. We regard it as defined over $\IQbar$; it is irreducible according to our convention."
  \textit{Source: Section 6 of \cite{MR4276287-mordell-lang-curves} (read from
  references/MR4276287-mordell-lang-curves-source/SettingUp.tex).}
  Fix integers \(g \ge 2\) and \(\ell \ge 3\). The \emph{Siegel moduli space}
  \(\mathbb{A}_g\) is the fine moduli space of principally polarized abelian
  varieties of dimension \(g\) with level-\(\ell\) structure. More precisely,
  \(\mathbb{A}_g\) represents the functor associating to each
  \(\mathrm{Spec}\,\mathbb{Z}[1/\ell]\)-scheme \(T\) the set of isomorphism
  classes of triples \((A \to T,\, \lambda,\, \alpha)\), where \(A \to T\) is an
  abelian scheme of relative dimension \(g\), \(\lambda\) is a principal
  polarization of \(A\), and \(\alpha\) is a level-\(\ell\) structure on \(A\).
  For \(\ell \ge 3\), this functor is representable, making \(\mathbb{A}_g\) a
  fine (not merely coarse) moduli scheme. Over \(\overline{\mathbb{Q}}\),
  \(\mathbb{A}_g\) is a smooth, irreducible, quasi-projective variety; see
  \cite[Theorem~7.9]{MFK:GIT94} or \cite[Theorem~1.9]{OortSteenbrink}.
\end{definition}

\begin{definition}[Torelli morphism]
  \label{def:torelli_map}
  \lean{MR4276287UniformityInMordellLangForCurves.Torelli.torelliMap}
  \uses{def:Mg_level, def:Ag_level, def:jacobian_av}
  % SOURCE: Section 6 of the paper (read from references/MR4276287-mordell-lang-curves-source/SettingUp.tex)
  % SOURCE QUOTE: "As $\mathbb{A}_g$ is a fine moduli space we have the following Cartesian diagram
  % \xymatrix{
  % \mathrm{Jac}(\mathfrak{C}_g) \ar[r] \ar[d]  \pullbackcorner & \mathfrak{A}_g \ar[d]^{\pi} \\
  % \mathbb{M}_g \ar[r]_{\tau} & \mathbb{A}_g
  % }
  % the bottom arrow is the Torelli morphism."
  \textit{Source: Section 6 of \cite{MR4276287-mordell-lang-curves} (read from
  references/MR4276287-mordell-lang-curves-source/SettingUp.tex).}
  The \emph{Torelli morphism} is the morphism of schemes
  \[
    \tau \colon \mathbb{M}_g \longrightarrow \mathbb{A}_g
  \]
  sending the moduli point \([C, \alpha]\) of a smooth genus-\(g\) curve \(C\)
  with level-\(\ell\) structure \(\alpha\) to the isomorphism class of its
  principally polarized Jacobian \([\mathrm{Jac}(C),\, \theta_C,\, \alpha']\),
  where \(\theta_C\) is the canonical principal polarization of \(\mathrm{Jac}(C)\)
  induced by the theta divisor, and \(\alpha'\) is the induced level-\(\ell\)
  structure.

  Equivalently, \(\tau\) is the bottom arrow in the Cartesian diagram
  \[
    \begin{array}{ccc}
      \mathrm{Jac}(\mathfrak{C}_g) & \longrightarrow & \mathfrak{A}_g \\
      \Big\downarrow & & \Big\downarrow^{\pi^{\mathrm{univ}}} \\
      \mathbb{M}_g & \underset{\tau}{\longrightarrow} & \mathbb{A}_g ,
    \end{array}
  \]
  where the left vertical arrow is the relative Jacobian of the universal curve
  \(\mathfrak{C}_g \to \mathbb{M}_g\) and the Cartesian square expresses that
  the relative Jacobian is obtained by base change from the universal abelian
  variety along \(\tau\).
\end{definition}

\begin{theorem}[Torelli's theorem: injectivity on geometric points]
  \label{thm:torelli_injective}
  \lean{MR4276287UniformityInMordellLangForCurves.Torelli.torelliInjective}
  \uses{def:torelli_map}
  \textit{Source: Classical Torelli theorem (no local reference file consulted for
  this standard result).}
  The Torelli morphism \(\tau \colon \mathbb{M}_g \to \mathbb{A}_g\) is
  injective on geometric points: if \(C\) and \(C'\) are smooth projective
  curves of genus \(g \ge 2\) over \(\overline{\mathbb{Q}}\) and there exists
  an isomorphism of principally polarized abelian varieties
  \((\mathrm{Jac}(C), \theta_C) \xrightarrow{\;\sim\;} (\mathrm{Jac}(C'), \theta_{C'})\),
  then \(C \cong C'\) as curves.
\end{theorem}

\begin{proof}
  \uses{def:torelli_map, def:jacobian_av}
  This is the classical Torelli theorem; see, e.g., \cite[Chapter~12]{BirkenhakeLange}
  or \cite[IV.1]{Mumford:CurveAV}. We recall the sketch.

  Given an isomorphism \(\varphi \colon (\mathrm{Jac}(C), \theta_C) \xrightarrow{\sim}
  (\mathrm{Jac}(C'), \theta_{C'})\) of principally polarized abelian varieties, the
  theta divisor \(\Theta_C \subset \mathrm{Jac}(C)\) maps to \(\Theta_{C'}\) (up to
  translation). The Riemann singularity theorem characterizes the locus of singular
  points of \(\Theta\) as the image of \(C^{(g-2)} \to \mathrm{Jac}(C)\), and from
  this one recovers the curve \(C\) (up to isomorphism) from the pair
  \((\mathrm{Jac}(C), \Theta_C)\): the variety \(W_{g-1} = \{[D] : D \text{ effective},
  \deg D = g-1\} \subset \mathrm{Jac}(C)\) (a translate of \(\Theta_C\)) has a
  natural \((g-1)\)-dimensional structure that encodes the curve.

  In the Lean formalization this theorem is taken as an axiom stub, to be verified
  against the appropriate Mathlib formalization once available.
\end{proof}

\begin{lemma}[Fiber of universal abelian variety over the Torelli image]
  \label{lem:universal_jacobian_fib}
  \lean{MR4276287UniformityInMordellLangForCurves.Torelli.universalJacobianFiber}
  \uses{def:torelli_map, def:universal_av, def:jacobian_av}
  % SOURCE: Section 6 of the paper (read from references/MR4276287-mordell-lang-curves-source/SettingUp.tex)
  % SOURCE QUOTE: "Let for the moment $S\rightarrow \mathbb{M}_g$ be the identity. If $s \in \mathbb{M}_g(\IQbar)$, then $\mathfrak{C}_s$ is the curve parametrized by $s$, and $\mathfrak{A}_{g,\tau(s)}$ is its Jacobian."
  \textit{Source: Section 6 of \cite{MR4276287-mordell-lang-curves} (read from
  references/MR4276287-mordell-lang-curves-source/SettingUp.tex).}
  For every \(s \in \mathbb{M}_g(\overline{\mathbb{Q}})\), let \(\mathfrak{C}_s\)
  denote the genus-\(g\) curve classified by \(s\). Then the fiber
  \[
    \mathfrak{A}_{g,\tau(s)} = (\pi^{\mathrm{univ}})^{-1}(\tau(s))
  \]
  of the universal abelian variety over the Torelli image \(\tau(s) \in
  \mathbb{A}_g(\overline{\mathbb{Q}})\) is canonically isomorphic (as a
  principally polarized abelian variety) to \(\mathrm{Jac}(\mathfrak{C}_s)\).
\end{lemma}

\begin{proof}
  \uses{def:torelli_map, def:universal_av, def:jacobian_av}
  This is a direct consequence of the Cartesian diagram defining \(\tau\)
  (see \cref{def:torelli_map}) and the fine moduli property of \(\mathbb{A}_g\)
  (see \cref{def:Ag_level}).

  By definition, \(\tau(s)\) is the moduli point classifying
  \((\mathrm{Jac}(\mathfrak{C}_s), \theta_{\mathfrak{C}_s})\). Since
  \(\mathbb{A}_g\) is a fine moduli space, pulling back the universal abelian
  variety \(\mathfrak{A}_g \to \mathbb{A}_g\) along the map
  \(\mathrm{Spec}\,\overline{\mathbb{Q}} \xrightarrow{\tau(s)} \mathbb{A}_g\)
  recovers the classified object. Thus
  \[
    \mathfrak{A}_g \times_{\mathbb{A}_g} \mathrm{Spec}\,\overline{\mathbb{Q}}
    \;\cong\; \mathrm{Jac}(\mathfrak{C}_s),
  \]
  which is precisely the statement that the fiber over \(\tau(s)\) is
  \(\mathrm{Jac}(\mathfrak{C}_s)\).

  In the Cartesian diagram of \cref{def:torelli_map}, the upper-left vertex is the
  relative Jacobian \(\mathrm{Jac}(\mathfrak{C}_g)\) over \(\mathbb{M}_g\), and the
  diagram is Cartesian: for each \(s \in \mathbb{M}_g(\overline{\mathbb{Q}})\),
  the fiber of \(\mathrm{Jac}(\mathfrak{C}_g)\) over \(s\) equals the fiber of
  \(\mathfrak{A}_g\) over \(\tau(s)\), i.e.,
  \(\mathrm{Jac}(\mathfrak{C}_s) \cong \mathfrak{A}_{g,\tau(s)}\).

  In the Lean formalization this is an axiom stub matching the fine moduli
  universal property.
\end{proof}

\section{Betti map and Betti form}

\begin{definition}[Universal principally polarized abelian variety]
  \label{def:universal_av}
  \lean{MR4276287UniformityInMordellLangForCurves.Torelli.universalAV}
  \uses{def:Ag_level}
  % SOURCE: Section 2 of the paper (read from references/MR4276287-mordell-lang-curves-source/BettiMapForm.tex)
  % SOURCE QUOTE: "We start to prove Proposition~\ref{PropBettiMap} for $S = \mathbb A_g$, the moduli space of principally polarized abelian variety of dimension $g$ with level-$\ell$-structure;  it is a fine moduli space. Let $\pi^{\mathrm{univ}} \colon \mathfrak{A}_g \rightarrow \mathbb A_g$ be the universal abelian variety."
  \textit{Source: Section 2 of \cite{MR4276287-mordell-lang-curves} (read from
  references/MR4276287-mordell-lang-curves-source/BettiMapForm.tex); also Section 6 (read
  from references/MR4276287-mordell-lang-curves-source/SettingUp.tex).}
  The \emph{universal principally polarized abelian variety} is the abelian scheme
  \[
    \pi^{\mathrm{univ}} \colon \mathfrak{A}_g \longrightarrow \mathbb{A}_g
  \]
  representing the universal object of the fine moduli functor of \cref{def:Ag_level}.
  For each geometric point \(s \in \mathbb{A}_g\), the fiber
  \(\mathfrak{A}_{g,s} = (\pi^{\mathrm{univ}})^{-1}(s)\) is the principally
  polarized abelian variety classified by \(s\). In particular, \(\pi^{\mathrm{univ}}\)
  is a principally polarized abelian scheme over \(\mathbb{A}_g\) equipped with
  a level-\(\ell\) structure, and it is projective over \(\mathbb{A}_g\); see
  \cite[Remark~\ref*{rmk:projembedding}]{MR4276287-mordell-lang-curves}.
\end{definition}
